\documentclass[11pt]{article}
\usepackage[a4paper,margin=27mm]{geometry}
\usepackage{amsmath,amssymb,amsthm,mathtools,booktabs,microtype}
\usepackage[T1]{fontenc}
\usepackage{lmodern}
\usepackage[colorlinks=true,linkcolor=blue,citecolor=blue,urlcolor=blue]{hyperref}
\usepackage{fancyhdr}
\pagestyle{fancy}\fancyhf{}
\lhead{Focused main theorem for review}\rhead{26 October 2026}\cfoot{\thepage}
\setlength{\headheight}{14pt}
\newtheorem{theorem}{Theorem}[section]
\newtheorem{lemma}[theorem]{Lemma}
\newtheorem{corollary}[theorem]{Corollary}
\newtheorem{proposition}[theorem]{Proposition}
\theoremstyle{definition}\newtheorem{definition}[theorem]{Definition}
\theoremstyle{remark}\newtheorem{remark}[theorem]{Remark}
\newcommand{\F}{\mathbb F}
\newcommand{\Mat}{\operatorname{Mat}}
\newcommand{\rank}{\operatorname{rank}}
\newcommand{\Span}{\operatorname{span}}
\newcommand{\Coeff}{\operatorname{Coeff}}
\newcommand{\im}{\operatorname{im}}
\newcommand{\SV}{\operatorname{SV}}
\newcommand{\supp}{\operatorname{supp}}
\title{The Interval-Rank Method\\for Unknown-Order ROABP Identity Testing}
\author{icekylinx \and GPT-6 Astra Pro}
\date{26 October 2026}
\begin{document}
\maketitle
\begin{center}
\fbox{\parbox{0.91\textwidth}{\textbf{Research status.} This document contains the complete proposed main argument and its geometric refinement, independently of all width-three classifications. It has internal model-assisted independent audits, but no external human expert review or proof-assistant certification. The larger checkpoint improves some constants and gives separate small-width results. This document is intended to make the core mathematical claim directly reviewable.}}
\end{center}
\begin{abstract}
We propose a deterministic black-box identity test for read-once oblivious
algebraic branching programs of every fixed width, with arbitrary unknown
variable order, arbitrary individual degree, and possibly singular factors.
The first step proves that products whose every contiguous interval
preserves its generic rank at zero have coefficient concentration bounded
only by their width. A descending induction on rank then constructs a
universal Shpilka--Volkovich shift preserving all interval ranks. The
elementary argument gives $2^{O(w^2)}$ seed pairs. Preserving scalar-saturated
image closures strengthens this to $2^{O(w\log w)}$ seed pairs. A rational
reparameterization yields $(nd+1)^{2^{O(w\log w)}}$ queries for individual
degree $d$, over any characteristic when an adequate evaluation field is
available. Over small finite fields, extension-field oracle access is an
explicit model requirement. This is a complete proposed proof, internally
audited but not externally peer reviewed; its validity and novelty remain
subject to that review.
\end{abstract}
\tableofcontents
\newpage
\section{Target, scope, and comparison}

A width-$w$ ROABP computes a polynomial of the form
\[
 f(x_1,\ldots,x_n)=uA_1(x_{\pi(1)})\cdots A_n(x_{\pi(n)})v,
\]
where $u\in\F^{1\times w}$, $v\in\F^{w\times1}$, and each $A_i$ is a
univariate polynomial matrix. The order $\pi$ is arbitrary and unavailable
to the black-box algorithm. Constant factors and unused variables are
allowed. Rectangular ROABPs of width at most $w$ can be padded to this form.

This target must be distinguished from a known variable order, and from
the stronger promise that a polynomial has a small ROABP in every order.
Polynomial hitting sets for constant width with a known order were given
by Gurjar, Korwar, and Saxena~\cite{GKS17}. The verified general unknown-order
baseline of Guo and Gurjar~\cite{GG20}, for fixed width, is
$n^{O(\log n/\log\log n)}\operatorname{poly}(d)$ under its stated field
conditions. Earlier work handles invertible-factor constant-width ROABPs
and arbitrary width two~\cite{AGKS15}.

Hu, van Melkebeek, and Morgan~\cite{HMM24} study the vanishing ideal of the
SV generator. Their ROABP theorem has an additional low-degree-monomial
condition. Theorem 1.9 and the end-of-Section-1 discussion explicitly identify full
black-box derandomization of arbitrary-degree constant-width ROABPs as
unresolved there. The proof below is intended to address that broader
target. A literature search can establish this comparison, but cannot by
itself certify that no later or differently phrased prior result exists.

A competing preprint of Singh and Singh~\cite{SS26}, version 1 of
13 February 2026, also claims an order-oblivious hitting set of polynomial
size for fixed width. Our separate research checkpoint records explicit
objections to intermediate statements in that version. These objections
do not certify the present proof or settle priority. The argument below
does not use that preprint, and both validity and any novelty claim require
independent scrutiny.

\section{Coefficient spaces and formal units}

All polynomial variable blocks in a product are pairwise disjoint. They
may contain several variables, or be empty. The support of a monomial
counts the blocks in which it has positive degree. When the blocks are
singletons, this is ordinary variable support.

\begin{definition}
For a matrix polynomial $M$, let $\Coeff(M)$ be the $\F$-linear span of its
coefficient matrices. We call $M$ \emph{$\ell$-concentrated} if this span is
generated by coefficients of block support at most $\ell$. The same
definition applies to row polynomials. Zero and constant polynomials
are $0$-concentrated.
\end{definition}

The concentration convention here is ``at most $\ell$''; some papers use
``strictly less than $\ell$''. This only changes indexing.

\begin{lemma}[Composition]\label{lem:composition}
Let $P(X)$ and $Q(Y)$ have disjoint blocks and concentration bounds $a,b$.
Then $P(X)Q(Y)$ is $(a+b)$-concentrated. Constant left and right matrices
preserve concentration. If a product is concentrated in at most $k$
macroblocks, and each macroblock is $b$-concentrated in its own original
blocks, the product is $kb$-concentrated in the original blocks.
\end{lemma}
\begin{proof}
Each coefficient of a product on disjoint blocks is a product of
coefficient matrices. Expand these matrices in their low-support spanning
sets and distribute. All resulting products are actual coefficients on
the disjoint original blocks. Constant transformations follow by linearity.
\end{proof}

\begin{lemma}[Formal units]\label{lem:units}
Let $H(X)$ be a matrix of formal power series, and let $R(X)$ be a square
formal matrix with invertible constant term. The linear spans of all
coefficient rows of $H$ and $RH$ coincide.
\end{lemma}
\begin{proof}
Every coefficient of $RH$ is a finite linear combination of coefficient
rows of $H$. The inverse $R^{-1}$ exists in the formal power series ring,
so applying the same argument to $H=R^{-1}(RH)$ gives the reverse inclusion.
\end{proof}

All generic ranks are over the rational-function field in the relevant
variables. Polynomial translation by a vector in an extended coefficient
field preserves these ranks, since it is an automorphism of that
polynomial ring. Specialization never increases generic rank.

\section{The normalized constant-rank lemma}

\begin{theorem}[Normalized row concentration]\label{thm:normalized}
Fix $P=\operatorname{diag}(I_r,0)\in\Mat_w(\F)$, with $r\ge1$. Let
$A_i(X_i)$ have generic rank $r$ and $A_i(0)=P$, on disjoint variable
blocks. For every constant row $u$, the row polynomial
$uA_1\cdots A_m$ is $2r$-concentrated.
\end{theorem}

\begin{proof}
Put $N_i=uA_1\cdots A_i$, $S_i=\Coeff(N_i)$, and let $T_i\subseteq\F^r$
be the projection of $S_i$ to its first $r$ coordinates. Setting the current
block to zero gives $N_i|_{X_i=0}=N_{i-1}P$, whence
\[
 T_{i-1}\subseteq T_i,\qquad T_{i-1}\times\{0\}\subseteq S_i.
\]
Disjointness also gives
\[
 S_i=\Span\{v\,\operatorname{coeff}_\alpha(A_i):v\in S_{i-1},\ \alpha\}.
\]

Write $E_i$ for the upper-left $r\times r$ block of $A_i$. Since
$E_i(0)=I_r$ and $\rank A_i=r$, Schur complementation over the rational
function field gives
\begin{equation}\label{eq:schur}
 A_i=\begin{bmatrix}I_r\\D_i\end{bmatrix}E_i
       \begin{bmatrix}I_r&C_i\end{bmatrix},
 \qquad E_i(0)=I_r,\quad C_i(0)=D_i(0)=0.
\end{equation}
All entries in this factorization are regular at zero, so the identity
holds in the formal power series ring as well.

We first prove a stable-step claim. Suppose $T_i=T_{i-1}=T$ and
$T\times\{0\}\subseteq S_{i-1}$. Then
\begin{equation}\label{eq:stable}
 S_i=\Span\Coeff\{(t,0)A_i:t\in T\}.
\end{equation}
For $v\in S_{i-1}$, every coefficient of the first $r$ coordinates of
$vA_i$ belongs to $T$. Equation~\eqref{eq:schur} writes $vA_i$ as that
active core times $[I_r,C_i]$. Thus its coefficients lie in the coefficient
span of $T[I_r,C_i]$.

For the reverse comparison, choose a basis row matrix $B$ of $T$. Since
$(T,0)\subseteq S_{i-1}$ and the active span is stable, $BE_i=R_iB$ for a
formal square matrix $R_i$ with $R_i(0)=I$. Lemma~\ref{lem:units} implies
that $BE_i[I_r,C_i]$ and $B[I_r,C_i]$ have identical coefficient-row spans.
The first family is contained in $S_i$. This proves~\eqref{eq:stable}.

Now partition the indices into maximal phases of constant $T_i$. At a
phase start $j$, suppose $S_j$ has a spanning set of actual coefficients
of support at most $L$. The next prefix is $(L+1)$-concentrated by direct
multiplication. At a later index $i\ge j+2$ in the same phase,
equation~\eqref{eq:stable} applies. Project the chosen prefix coefficients
at $j$ to $T$, set blocks $j+1,\ldots,i-1$ to zero, and activate block $i$.
There is at least one intervening $P$, so these are exactly coefficients
realizing $(t,0)A_i$, and their support is at most $L+1$. A long stable
phase therefore incurs only one additional block.

A useful invariant is: at the start of a phase of active dimension $t$,
support at most $2t-1$ spans the row space; within the phase, $2t$ suffices.
It holds initially since $S_1$ is $1$-concentrated. If $N_1\ne0$, then
$T_1\ne0$ by~\eqref{eq:schur}. A growth step from dimension $t$ needs at
most $2t+1\le2\dim T_{\mathrm{new}}-1$, establishing the invariant.
Since $\dim T_i\le r$, the final bound is $2r$. If $N_1=0$, the entire
product is zero and the assertion is immediate.
\end{proof}

\begin{corollary}[Full matrix concentration]\label{cor:matrix}
Under the hypotheses of Theorem~\ref{thm:normalized}, the full matrix
product $A_1\cdots A_m$ is $2wr$-concentrated.
\end{corollary}
\begin{proof}
Use the row $[e_1^T,\ldots,e_w^T]$ of length $w^2$ and factors
$I_w\otimes A_i$. Their common constant term is a projector of rank $wr$,
conjugate by a fixed permutation to a standard projector. The output row
is the row-wise vectorization of the full matrix product. Apply
Theorem~\ref{thm:normalized} with rank $wr$.
\end{proof}

The distinction between rows and full matrices is essential: separately
proving a $2r$ bound for each row does not establish a $2r$ bound for the
joint matrix coefficient space.

\begin{proposition}[Sharpness of the row support bound]
For every $r\ge1$, width $r+1$ admits an example satisfying the normalized
hypotheses whose scalar output is a monomial of support $2r$.
\end{proposition}
\begin{proof}
Let $q=r+1$ and $P=\operatorname{diag}(I_r,0)$. Take, for $j=1,\ldots,r$,
the consecutive factors
\[
 P+x_{2j-1}E_{qj},\qquad P+x_{2j}E_{jq}.
\]
Each has generic rank $r$ and constant term $P$. In the scalar contraction
from row $e_q^T$ to column $e_q$, the only surviving path alternates through
the active channels $1,\ldots,r$ and the inactive channel $q$. Its weight
is $x_1\cdots x_{2r}$. Any omitted transition either leaves the inactive
channel across a zero diagonal entry or strands the path in an active
channel that cannot reach the final inactive endpoint.
\end{proof}

\section{Rank profiles imply concentration}

\begin{lemma}[Constant gauges]\label{lem:gauge}
Suppose constant matrices $B_1,\ldots,B_k$ have rank $r$ and every adjacent
product $B_iB_{i+1}$ has rank $r$. There are invertible matrices
$T_0,\ldots,T_k$ such that
\[
 T_{i-1}^{-1}B_iT_i=\operatorname{diag}(I_r,0).
\]
In particular, the full product has rank $r$.
\end{lemma}
\begin{proof}
Choose an ordered basis for a complement to $\ker B_k$ at the last cut.
Propagate these active basis vectors backwards through the maps $B_i$.
At an internal cut the active space is $\im B_{i+1}$, transverse to
$\ker B_i$ because $\rank(B_iB_{i+1})=r$. Complete each active basis by a
basis of $\ker B_i$. Complete the left endpoint basis arbitrarily. In these
bases every map is the stated projector. Their product is the same
projector.
\end{proof}

For a positive integer $q$ and an allowed-rank set
$\mathcal R\subseteq\{1,\ldots,q\}$ containing $q$, define constants by
descending rank:
\begin{align}
 a_{q,\mathcal R}(q)&=q^2-1,\label{eq:base}\tag{4.1}\\
 a_{q,\mathcal R}(s)&=2qs\left(2\max_{t\in\mathcal R,\ t>s}
              a_{q,\mathcal R}(t)+1\right),\quad s<q.
 \label{eq:recurrence}\tag{4.2}
\end{align}

\begin{theorem}[Rank-profile concentration]\label{thm:profile}
Let $M=A_1(X_1)\cdots A_m(X_m)$ have width $q$. Suppose every contiguous
interval product has the same generic rank as its constant specialization,
and every nonzero interval rank belongs to $\mathcal R$. If $M$ has rank
$s>0$, it is $a_{q,\mathcal R}(s)$-concentrated. If $M$ has rank zero,
it is the zero matrix polynomial.
\end{theorem}
\begin{proof}
Induct downwards through $\mathcal R$ on the rank of the full product.
If $s=q$, every factor has invertible constant term. Constant gauges make
all these terms the identity. The prefix matrix coefficient spaces are
nested subspaces of $\Mat_q(\F)$, initially spanned by $I_q$. At a strict
dimension increase, multiplying a prior coefficient spanning set adds at
most one active block. At a stable step the old coefficients remain actual
coefficients by choosing the identity constant. There are at most $q^2-1$
increases, proving~\eqref{eq:base}.

Suppose $s<q$. Starting from the left, take a shortest prefix of rank $s$,
then repeat if the remaining product still has rank $s$. Write the chosen
blocks as $B_j=C_jA_{e_j}$, where $C_j$ is the prefix before the final
original factor. Its rank is greater than $s$, unless it is empty, in
which case it is the identity. The remaining tail $T$ also has rank
greater than $s$, or is empty. This process is valid because the rank of
any subinterval is at least the rank of the full product.

Every interval inside $C_j$ or $T$ has rank greater than $s$. By induction,
these products are concentrated at
$b=\max_{t\in\mathcal R,t>s}a_{q,\mathcal R}(t)$.
Absorb $T$ into the final $B_j$. Each resulting macroblock is
$(2b+1)$-concentrated and has both generic and constant rank $s$.
Every adjacent macroblock product has rank $s$: it contains a rank-$s$
factor and is an interval of a full rank-$s$ product. The constant rank is
also $s$ by hypothesis.

Lemma~\ref{lem:gauge} simultaneously normalizes the macroblock constants.
Corollary~\ref{cor:matrix} concentrates their product in at most $2qs$
macroblocks. Lemma~\ref{lem:composition} replaces each macroblock coefficient
by original coefficients, giving $2qs(2b+1)$ and proving the recurrence.
\end{proof}

\section{A universal rank-preserving SV shift}

Choose distinct $a_1,\ldots,a_n\in\F$ and let
\[
 L_i(y)=\prod_{j\ne i}\frac{y-a_j}{a_i-a_j},\qquad
 \SV_{\ell,i}(y,z)=\sum_{j=1}^{\ell}z_jL_i(y_j).
\]
The generator is additive in disjoint seed pairs:
$\SV_a+\SV_b=\SV_{a+b}$ after relabeling. It includes the zero vector.

\begin{lemma}[Support hitting]\label{lem:sv}
Over any coefficient-field extension of $\F$, a polynomial with a nonzero
monomial of support at most $\ell$ remains nonzero under $\SV_\ell$.
\end{lemma}
\begin{proof}
Set the corresponding $y_j$ to the abscissas of its support, and set unused
$z_j$ to zero. This leaves those coordinates as independent variables and
sets all other coordinates to zero. The chosen coefficient remains nonzero.
\end{proof}

For $1\le r<w$, put $W_r=\binom wr$ and
\[
 \mathcal R_r=\left\{\binom sr:r+1\le s\le w\right\},\qquad
 M_r=\max_{s\in\mathcal R_r}a_{W_r,\mathcal R_r}(s),\qquad
 h_r=2M_r+2.
\]
Finally define
\begin{equation}\label{eq:k}
 k(w)=1+\sum_{r=1}^{w-1}h_r,\qquad k(1)=1.
\end{equation}

\begin{theorem}[Universal interval rank preservation]\label{thm:preserve}
For every product of width-$w$ univariate polynomial matrices on distinct
variables, in any order, $\SV_{k(w)}$ preserves the generic rank of every
contiguous interval, over the rational-function field in its seeds.
The same assertion allows constant factors without variables.
\end{theorem}
\begin{proof}
Construct $G_w,G_{w-1},\ldots,G_1$ so that $G_r$ preserves all interval ranks
at least $r$. Take $G_w=\SV_1$. A full-rank interval consists of invertible
univariate factors and invertible constant factors. A nonzero univariate
determinant stays nonzero under its coordinate $z_1L_i(y_1)$, so this base
case holds over $\F(y_1,z_1)$.

Suppose $G_{r+1}$ has been constructed. Work over its seed field $K$ and
translate the original variables by $G_{r+1}$. The residual variables
still occur in disjoint layers. Fix any interval of generic rank exactly
$r$. Partition it greedily as
\[
 B_1\cdots B_qT,\qquad B_j=C_jA_{e_j},
\]
where each $B_j$ is a shortest rank-$r$ prefix and $T$ has rank greater
than $r$, or is empty. Every interval inside each $C_j$ and inside $T$
has rank greater than $r$, hence that rank is preserved by the previous
shift at residual zero.

Take the $r$th exterior power. Its dimension is $W_r$, multiplication is
preserved, and the exterior power of a rank-$s$ matrix has rank
$\binom sr$ over every field. Consequently the exterior interval products
inside $C_j$ or $T$ satisfy the hypotheses of
Theorem~\ref{thm:profile} with allowed ranks $\mathcal R_r$. Each of
$\bigwedge^r C_j$ and $\bigwedge^r T$ is $M_r$-concentrated. A single
factor $\bigwedge^r A_{e_j}$ is univariate and thus $1$-concentrated.

Lemma~\ref{lem:composition} gives concentration bounds
\[
\begin{array}{c|ccc}
\text{product}&\bigwedge^r B_j&\bigwedge^r(B_jB_{j+1})&\bigwedge^r(B_qT)\\
\hline
\text{bound}&M_r+1&2M_r+2&2M_r+1.
\end{array}
\]
Each underlying product has generic rank $r$. Some exterior entry is
therefore nonzero, and concentration gives such an entry with a nonzero
coefficient of support at most $h_r$. Lemma~\ref{lem:sv} shows that it
survives a fresh $\SV_{h_r}$. Thus, at the new specialization, each $B_j$
has rank $r$, each adjacent pair has rank $r$, and $B_qT$ has rank $r$.

These statements imply that the entire interval still has rank $r$.
For example, write $B_j=U_jV_j$ with $U_j,V_j$ of full column and row
rank $r$. All transfers $V_jU_{j+1}$ are invertible, and $V_qT$ has full
row rank. Their product has rank $r$.

Set $G_r=G_{r+1}+\SV_{h_r}$. Already preserved minors of higher rank
cannot become identically zero, since setting the new $z$ seeds to zero
recovers the previous map. Specialization cannot raise rank. This proves
the induction. The same universal map works for each interval and each
ROABP; the interval-dependent greedy partition is only a proof device.
Additivity of SV gives $G_1=\SV_{k(w)}$.
\end{proof}

\begin{remark}[The tail is not absorbed in this proof]
The concentration proof in Section 4 may absorb its tail because it
already assumes preservation of every interval rank. In the construction
above, doing so would require a bound for $B_{q-1}B_qT$ with three
higher-rank runs. Instead we preserve $B_qT$ separately. Its image equals
$\im B_q$, and injectivity of the preceding transfers gives the required
full rank. This distinction prevents a missing term in the support budget.
\end{remark}

\section{The proposed hitting set and its size}

\begin{theorem}[Candidate main theorem]\label{thm:main}
Let $f$ be a nonzero $n$-variate polynomial of individual degree at most
$d$, computed by a width-at-most-$w$ ROABP in an arbitrary unknown order.
Then
\[
 f\bigl(\SV_{k(w)}\bigr)\not\equiv0.
\]
For $w\ge2$, $k(w)\le2^{4w^2}$ is a valid coarse bound. If the field
contains a set $S$ of more than $n^2d$ elements and the $n$ abscissas,
evaluating all $2k(w)$ seeds on $S$ produces a hitting set of size at most
\[
 |S|^{2k(w)}=(n^2d+1)^{2^{O(w^2)}}
\]
when a set of the indicated size is available. This is polynomial in
$n,d$ for every fixed $w$.
\end{theorem}
\begin{proof}
Pad the scalar product $uA_1\cdots A_nv$ by constant matrices
$e_1u$ and $ve_1^T$. The resulting full product is $fE_{11}$, of generic
rank one. Theorem~\ref{thm:preserve} preserves this rank and therefore
the nonzeroness of $f$.

For the quantitative bound, an allowed-rank recurrence at width $W_r$
has at most $w-r$ levels. Writing $a+1$ avoids the additive term and gives
\[
 M_r+1\le(5W_r^2)^{w-r}\le(5\cdot4^w)^{w-r}.
\]
Therefore $k(w)\le1+2\sum_{r=1}^{w-1}(5\cdot4^w)^{w-r}\le2^{4w^2}$
for $w\ge2$. The exact recurrence is much smaller than this coarse bound.

Each coordinate of SV has total seed degree at most $n$, while $f$ has
total degree at most $nd$. Thus $f(\SV)$ has degree at most $n^2d$.
The elementary grid lemma for nonzero polynomials supplies a nonzero
evaluation on $S^{2k(w)}$.
\end{proof}

The field issue is substantive: a formal polynomial such as $x^q-x$ is
nonzero over $\F_q$ yet vanishes at all base-field points. Extension-field
queries must therefore be explicitly allowed if the base field is too
small. No statement here converts a base-field-only oracle into an
extension-field oracle. Over characteristic zero, the construction can
use ordinary distinct rational abscissas and seed values.

\section{A geometric refinement: \texorpdfstring{$2^{O(w\log w)}$}{2 to O(w log w)} seed pairs}
\label{sec:geometry}

The preceding proof uses exterior powers to turn rank conditions into
scalar nonvanishing. Its loss comes from applying a linear coefficient-space
lemma at the exterior dimension. We can avoid this loss by preserving a
geometric object that controls every homogeneous polynomial in the matrix
entries simultaneously. This section gives the stronger candidate bound.
The earlier proof is retained as an independent, purely linear-algebraic
route to the qualitative constant-width conclusion.

Work over an algebraic closure $k$ of the current coefficient field. For a
polynomial matrix map $M(X)$, define its scalar-saturated image closure
\[
 \mathcal C(M)=\overline{\{\lambda M(x):\lambda\in k,\ x\in k^{|X|}\}}.
\]
The closure is Zariski closure. In particular, over a finite coefficient
field we do \emph{not} take only its finite set of rational points. The
variety $\mathcal C(M)$ is an irreducible affine cone, being the closure of
the image of an affine space under a polynomial map. The auxiliary scalar
$\lambda$ is not an original variable block and incurs no support charge.

Call a restriction $X_j=0$ outside a subset $J$ \emph{cone preserving} if
it leaves $\mathcal C(M)$ unchanged. The selected $J$ may depend on the
matrix product; no black-box algorithm must find it.

\begin{lemma}[Consequences of a cone-preserving restriction]
\label{lem:cone-composition}
If $J$ is cone preserving for $M$, every nonzero homogeneous polynomial
function $p(M)$ remains nonzero after that restriction. In particular,
every exterior-power matrix $\bigwedge^r M$ is $|J|$-concentrated in the
original blocks. Cone-preserving restrictions compose across disjoint
factors and across disjoint macroblocks, with the same support addition
and multiplication rules as Lemma~\ref{lem:composition}.
\end{lemma}
\begin{proof}
A homogeneous polynomial vanishes on an image if and only if it vanishes
on its scalar-saturated closure. This proves the first assertion. A linear
functional on the entries of $\bigwedge^r M$ is a homogeneous polynomial
of degree $r$ in the entries of $M$. If it vanishes on all coefficients
of the restriction, it vanishes on the full matrix polynomial. Duality
of finite-dimensional vector spaces therefore gives the coefficient-span
assertion.

For composition, products of dense subsets are dense in the corresponding
products of varieties, and multiplication of matrices is a polynomial map.
Independent scalar amplitudes in the factors multiply to one scalar
amplitude of the full product. Thus replacing each factor image by a
cone-dense restricted image preserves the product cone. The same reasoning
applies after unselected macroblocks have been fixed at zero.
\end{proof}

\begin{lemma}[Invertible constants]\label{lem:cone-invertible}
A width-$w$ product with invertible constant factors has a cone-preserving
restriction to at most $w^2-1$ blocks.
\end{lemma}
\begin{proof}
Normalize the constants to $I_w$ and let $M_i=A_1\cdots A_i$, $M_0=I_w$.
The cones $\mathcal C(M_i)$ are nested, since the current block may be set
to zero. They are irreducible, live in a space of dimension $w^2$, and
start at the line $kI_w$, of dimension one. A proper inclusion of
irreducible closed varieties increases dimension. At a stable step keep
the previous restriction. At a growth step retain the new block as well;
the multiplication morphism and density show that this suffices. There
are at most $w^2-1$ growth steps.
\end{proof}

\begin{lemma}[Normalized constant-rank cones]\label{lem:cone-normalized}
Under the hypotheses of Theorem~\ref{thm:normalized}, the full matrix
product has a cone-preserving restriction to at most $2wr$ blocks.
\end{lemma}
\begin{proof}
Put $M_i=A_1\cdots A_i$, let $Q_i$ be its first $r$ columns, and set
$\mathcal S_i=\mathcal C(M_i)$, $V_i=\mathcal C(Q_i)$. For $i\ge2$,
\[
 V_{i-1}\subseteq V_i,\qquad [V_{i-1},0]\subseteq\mathcal S_i.
\]
The $V_i$ are irreducible cones in $\Mat_{w\times r}$, so their dimensions
are at most $wr$. They have dimension at least one because $Q_i(0)$ is
the nonzero first-$r$-column part of $P$.

We prove the geometric stable-step identity. Suppose
$V_i=V_{i-1}=V$ and $[V,0]\subseteq\mathcal S_{i-1}$. Then
\begin{equation}\label{eq:geometric-stable}
 \mathcal S_i=
 \overline{\{[T,0]A_i(a):T\in V,\ a\in k^{|X_i|}\}}.
\end{equation}
Use the Schur factorization~\eqref{eq:schur} on the nonempty open set
$U$ where $E_i$ is invertible. For every $H\in\mathcal S_{i-1}$ and
$a\in U$, the active core of $HA_i(a)$ belongs to $V$, by stability
and continuity from the actual prefix image. The output equals this
core times $[I,C_i(a)]$.

Conversely, baseline inclusion implies $VE_i(a)\subseteq V$ for each
fixed $a\in U$. Right multiplication by $E_i(a)$ is an invertible
linear map, so $VE_i(a)$ is a closed irreducible variety of the same
dimension as $V$. Inclusion forces equality. Hence the baseline outputs
at this fixed $a$ are precisely
\[
 \{TE_i(a)[I,C_i(a)]:T\in V\}
   =\{T[I,C_i(a)]:T\in V\}.
\]
This gives both inclusions in~\eqref{eq:geometric-stable} after taking
closures. Restricting the current parameters to $U$ loses no image closure:
it is a dense open subset, and the full parameter map is polynomial.

Partition the indices into phases where $V_i$ is constant. At a phase
start $j$, retain a cone-preserving restriction of $M_j$ using $L$ blocks.
The next step costs at most $L+1$. At a later step $i\ge j+2$, project
that restricted image to its active core, set all intervening blocks to
zero, and retain block $i$. An intervening $P$ gives the baseline in
\eqref{eq:geometric-stable}, so this restriction is cone preserving and
still costs at most $L+1$.

A strict active-cone growth raises the next phase-start budget by at
most two. The first product costs one block, and there are at most
$wr-1$ later dimension increases. Thus $1+2(wr-1)+1=2wr$ blocks suffice.
\end{proof}

\begin{lemma}[Intrinsic-rank bound]\label{lem:intrinsic-cone}
The normalized cone restriction in Lemma~\ref{lem:cone-normalized} can
be chosen using at most $2r^2$ blocks, independently of the ambient width.
\end{lemma}
\begin{proof}
First consider a fixed $a\times w$ initial matrix $R$ with nonzero active
part, followed by normalized rank-$r$ factors. The preceding phase proof
applies to rectangular output matrices. Its active cones lie in
$\Mat_{a\times r}$ and start at the nonzero line spanned by the active
part of $R$, with support budget zero. There are at most $ar-1$ strict
growths. Paying at most two per growth and one for the final stable
phase gives $2ar-1$ blocks. The first stable step is paid separately,
so arbitrary passive columns of $R$ cause no difficulty.

Now retain the entire first original block $X_1$ and work over
$K=k(X_1)$. The first $r$ columns $U_1$ of $A_1$ have full column rank,
and the rank hypothesis gives $A_1=U_1R_1$, where $R_1=[I_r,C_1]$ is
a fixed $r\times w$ matrix over $K$. The rectangular statement selects
at most $2r^2-1$ of the remaining blocks while preserving the output cone
over an algebraic closure of $K$. Left multiplication by $U_1$ preserves
this equality. Including $X_1$ gives $2r^2$ blocks.

For completeness, descent to $k$ follows from homogeneous identities.
If a homogeneous polynomial in the full matrix entries vanishes on the
selected restriction, it vanishes there over $K$, hence on the unrestricted
cone over $K$, and hence as a polynomial identity in all original
variables over $k$. The subset is common to every such polynomial.
Homogeneous vanishing ideals determine closed cones, proving the claim.
\end{proof}

Define constants at the original width by
\begin{align*}
 D_w(w)&=w^2-1,\\
 D_w(r)&=2r^2\left(2\max_{r<s\le w}D_w(s)+1\right),\qquad 1\le r<w.
\end{align*}
\begin{theorem}[Rank-profile cone preservation]\label{thm:cone-profile}
If every contiguous interval of a width-$w$ product preserves its generic
rank at zero, a nonzero full product of rank $r$ has a cone-preserving
restriction to at most $D_w(r)$ blocks.
\end{theorem}
\begin{proof}
Use the same downward rank induction and the same greedy blocks
$B_j=C_jA_{e_j}$ as in Theorem~\ref{thm:profile}. The full-rank case is
Lemma~\ref{lem:cone-invertible}. The higher-rank prefixes and tail have
cone-preserving restrictions by induction. Absorbing the tail into the
last macroblock costs at most $2\max_{s>r}D_w(s)+1$ original blocks
per macroblock. Constant gauges and Lemma~\ref{lem:intrinsic-cone}
select at most $2r^2$ macroblocks. Lemma~\ref{lem:cone-composition} then
gives exactly the displayed recurrence.
\end{proof}

\begin{theorem}[Improved candidate generator]\label{thm:improved}
Let
\[
 k_{\mathrm{geo}}(w)=1+\sum_{r=1}^{w-1}
       \left(2\max_{r<s\le w}D_w(s)+2\right).
\]
Then $\SV_{k_{\mathrm{geo}}(w)}$ preserves every interval rank of every
unknown-order width-$w$ ROABP matrix product and hits every nonzero scalar
ROABP of that width. Moreover
\[
 k_{\mathrm{geo}}(w)\le4^{w-1}(w!)^2=2^{O(w\log w)}\quad(w\ge2),
 \qquad |\mathcal H|\le(n^2d+1)^{2^{O(w\log w)}}.
\]
The variant $t+\SV_{k_{\mathrm{geo}}(w)-1}$ uses one fewer seed variable.
\end{theorem}
\begin{proof}
Repeat the descending-rank construction of Theorem~\ref{thm:preserve}.
At rank $r$, every interval inside the higher-rank runs $C_j$ and tail
$T$ has already had its rank preserved. Theorem~\ref{thm:cone-profile}
and Lemma~\ref{lem:cone-composition} concentrate \emph{all} their
$r$th exterior entries at the original-width level
$\max_{s>r}D_w(s)$. The local pair and tail bounds are therefore
$2\max_{s>r}D_w(s)+2$, with no exterior-dimension substitution.
The same block-chain argument completes the induction.

For $w\ge2$, the $D_w(r)$ increase as $r$ decreases, and
\[
 D_w(r)+1\le4r^2\bigl(D_w(r+1)+1\bigr).
\]
Consequently
\[
 D_w(s)+1\le4^{w-s}\left(\frac{(w-1)!}{(s-1)!}\right)^2w^2.
\]
Summing from $s=2$ to $w$ and using a geometric upper bound gives
\[
 k_{\mathrm{geo}}(w)\le1+\frac83\,4^{w-2}(w!)^2
                  \le4^{w-1}(w!)^2\quad(w\ge2).
\]
The finite-grid argument
uses the original polynomial $f$, of degree at most $nd$, exactly as
before. All intermediate seeds remain symbolic; their ranks are not
preserved by prematurely fixing a numerical seed assignment.
\end{proof}

This geometric strengthening uses only elementary facts about irreducible
image closures and dimension, in addition to the previous matrix
factorizations. Its crucial new assertion is preservation of the entire
image cone under one small coordinate restriction. It is stronger than
linear coefficient concentration and is what justifies applying the same
width bound to minors of arbitrary order.


\section{Rational evaluations and the explicit query bounds}
\label{sec:rational}

The polynomial form of SV has degree $n-1$ in each interpolation seed.
A rational reparameterization removes this extra factor from the grid
degree. This is also a direct connection with the rational-function
evaluation framework of Hu, van Melkebeek, and Morgan~\cite{HMM24}.

Write $P(T)=\prod_i(T-a_i)$ and $d_i=P'(a_i)\ne0$. Then
\[
 L_i(y)=\frac{P(y)}{d_i(y-a_i)}.
\]
All uses of the initial common shift $t$ remain valid if coordinate $i$
is shifted by $c_it$ for arbitrary nonzero constants $c_i$: a nonzero
univariate polynomial remains nonzero at $c_it$. Thus the proved variants
also hold for $c_it+\SV_m$.

Apply this to the ROABP polynomial $g(x)=f(d_1x_1,\ldots,d_nx_n)$ with
$c_i=1/d_i$. Reparameterize $z_jP(y_j)$ as an independent residue
$\beta_j$; this is an invertible rational change of parameters over
$\F(y)$. We obtain the unweighted rational generator
\begin{equation}\label{eq:rational-generator}
 H_i(t,\beta,y)=t+\sum_{j=1}^{m}\frac{\beta_j}{y_j-a_i}.
\end{equation}
Whenever $c_it+\SV_m$ has the preceding hitting property, so does $H$.
Its values are evaluations, at the abscissas $a_i$, of a univariate
rational function with numerator and denominator degrees at most $m$.

\begin{lemma}[Grid degree after clearing poles]\label{lem:rational-grid}
Let $f$ have individual degree at most $d$ and suppose $f(H)\not\equiv0$.
Then
\[
 N(t,\beta,y)=
   \left(\prod_{i=1}^n\prod_{j=1}^{m}(y_j-a_i)^d\right)f(H)
\]
is a nonzero polynomial of degree at most $nd$ in \emph{each} of its
$2m+1$ seed variables. If $S,Y\subseteq\F$ have size $nd+1$ and
$Y\cap\{a_1,\ldots,a_n\}=\varnothing$, the evaluations on
$S^{m+1}\times Y^m$ are all defined and include a nonzero value of $f(H)$.
\end{lemma}
\begin{proof}
The denominator-clearing product is nonzero, so $N$ is nonzero. For a
fixed pole variable $y_j$, each coordinate has the form
$A_i+\beta_j/(y_j-a_i)$ with $A_i$ independent of $y_j$. A monomial
$\prod_i x_i^{e_i}$ of $f$, with $0\le e_i\le d$, becomes after clearing
these denominators a polynomial of $y_j$-degree at most
\[
 \sum_i e_i+\sum_i(d-e_i)=nd.
\]
The denominators involving other poles do not change this degree. The
degrees in $t$ and in each residue are at most $\sum_i e_i\le nd$.
The individual-degree grid lemma gives a nonzero value of $N$. All
denominators are nonzero on the chosen grid, so this is a nonzero value
of $f(H)$.
\end{proof}


\begin{corollary}[Rational finite grid for the elementary construction]
Let $k(w)$ be defined in~\eqref{eq:k}. For the class in Theorem~\ref{thm:main}, the map
\[
 H_i=t+\sum_{j=1}^{k(w)-1}\frac{\beta_j}{y_j-a_i}
\]
is a hitting map. If the evaluation field has at least $nd+n+1$ elements, enumerating $nd+1$ values for every seed, with pole seeds avoiding all abscissas, gives at most $(nd+1)^{2k(w)-1}$ queries.
\end{corollary}
\begin{proof}
The initial $\SV_1$ in Theorem~\ref{thm:preserve} is used only to preserve nonzero univariate determinants. It can therefore be replaced by the weighted common shift $(c_it)_i$ for any nonzero constants $c_i$. Every subsequent stage of that proof is unchanged. The weighted-shift rescaling above and Lemma~\ref{lem:rational-grid} give the assertion.
\end{proof}
Over $\mathbb Q$ one may take $a_i=i-1$, the amplitude and common-shift grid $\{0,\ldots,nd\}$, and the pole grid $\{n,\ldots,n+nd\}$. Query coordinates have bit length $O(k(w)\log(nd+n+1))$. The enumeration does not use the unknown variable order, the matrices, their ranks, or their coefficients.


\begin{corollary}[The improved finite-grid bound]\label{cor:focused-main}
Let $f\ne0$ have individual degree at most $d$ and have a width-at-most-$w$
ROABP in some variable order, which need not be given. Put
$K=k_{\mathrm{geo}}(w)$, with $K=1$ for $w=1$. Then
\[
 f\left(t+\sum_{j=1}^{K-1}\frac{\beta_j}{y_j-a_1},\ldots,
         t+\sum_{j=1}^{K-1}\frac{\beta_j}{y_j-a_n}\right)\ne0
\]
as a rational function of independent seeds. If the evaluation field has
at least $nd+n+1$ elements, the indicated grids give at most
\[
 (nd+1)^{2K-1},\qquad K\le4^{w-1}(w!)^2=2^{O(w\log w)}
 \quad(w\ge2)
\]
oracle queries. This is polynomial in $n,d$ for each fixed $w$. It is an
XP-type bound in the width parameter, not a claim of fixed-parameter
tractability with a width-independent polynomial exponent.
\end{corollary}
\begin{proof}
The only use of the first seed pair in Theorem~\ref{thm:improved} is the
full-rank base case of Theorem~\ref{thm:preserve}. Replace it by the weighted
common shift as above, leaving every later rank-descent step unchanged.
Apply the rational rescaling and Lemma~\ref{lem:rational-grid}.
\end{proof}

\section{Two direct algebraic consequences}
These consequences require no additional classification or algorithmic
closure property. They are recorded as implications of the proposed main
theorem, without a separate novelty claim.

\begin{corollary}[Factors and products of factors]\label{cor:factors}
The formal rational map in Corollary~\ref{cor:focused-main} is nonzero on
every nonzero polynomial factor of a nonzero width-at-most-$w$ ROABP.
It is also nonzero on any product of such factors, even if the parent
ROABPs use different variable orders. For a factor or product of known
individual degree at most $d'$, Lemma~\ref{lem:rational-grid} gives the
corresponding finite grid with $nd'+1$ values per seed.
\end{corollary}
\begin{proof}
Substitution by the rational map is an algebra homomorphism from
$\F[X]$ into the field of rational functions in the seeds. If $F=gh$ is
a nonzero parent polynomial, then $F(H)=g(H)h(H)\ne0$ by the main theorem,
so both factors remain nonzero. Products of nonzero rational functions
remain nonzero. No bound on the ROABP width of a factor is asserted or
needed. The finite grid uses the degree of the polynomial being tested,
not a degree bound silently inherited from an unspecified parent.
\end{proof}

\begin{corollary}[A degree-unrestricted kernel-ideal lower bound]
Fix distinct abscissas and let $I_m$ be the kernel of substitution
$x_i\mapsto t+\sum_{j=1}^m\beta_j/(y_j-a_i)$ into the rational-function
field. If $m\ge k_{\mathrm{geo}}(w)-1$, then $I_m$ contains no nonzero
polynomial with a width-at-most-$w$ ROABP in any variable order, regardless
of its individual degree. Consequently every nonzero $f\in I_m$ has
ROABP width
\[
 \Omega\!\left(\frac{\log(m+2)}{\log\log(m+3)}\right)
\]
in every order, as $m$ grows. When $n>2m+1$, the ideal is nonzero.
\end{corollary}
\begin{proof}
If more poles are available than required by the main theorem, set their
residues to zero. A nonzero composition on that specialization cannot
be identically zero before specialization. Thus a nonzero width-$w$
element in the kernel would force $m<k_{\mathrm{geo}}(w)-1$.
The upper bound $k_{\mathrm{geo}}(w)\le4^{w-1}(w!)^2$ implies
$\log(m+2)=O(w\log(w+1))$, giving the stated inversion. Finally the target
field has transcendence degree $2m+1$; if there are more source variables,
their images are algebraically dependent. This last assertion does not
provide an explicit kernel element or strengthen the width lower bound.
\end{proof}

\section{What an external reviewer needs to check}
The mathematical dependencies are fully contained above. The crucial checks are: the stable-step identity in the normalized row lemma; joint full-matrix concentration through vectorization; simultaneous constant gauges; rank-profile induction on shortest rank-drop prefixes; and the three separate local witnesses in the rank descent, including its final tail. At every descent stage old seeds are retained as coefficient-field parameters. No numerical choice of those seeds occurs before the final grid. The variable-order-dependent partition is used only inside the correctness proof.

The elementary qualitative claim does not depend on the geometric section. The improved width dependence does use that section, especially its stable-step equality and its descent from a generic first block to a single coordinate restriction over the original field. Neither argument depends on the weighted sharp-row theorem, any width-three invariant classification, or the separate critique of a competing preprint. Those separate results are not certified by checking this manuscript.
\begin{thebibliography}{9}
\bibitem{AGKS15} M. Agrawal, R. Gurjar, A. Korwar, and N. Saxena.
\emph{Hitting-sets for ROABP and sum of set-multilinear circuits}.
SIAM Journal on Computing, 2015. Preprint:
\url{https://arxiv.org/abs/1406.7535}.
\bibitem{GKS17} R. Gurjar, A. Korwar, and N. Saxena.
\emph{Identity Testing for Constant-Width, and Any-Order, Read-Once
Oblivious Arithmetic Branching Programs}. Theory of Computing 13(2), 2017.
\url{https://theoryofcomputing.org/articles/v013a002/}.
\bibitem{GG20} Z. Guo and R. Gurjar.
\emph{Improved Explicit Hitting-Sets for ROABPs}. APPROX/RANDOM 2020.
\url{https://doi.org/10.4230/LIPIcs.APPROX/RANDOM.2020.4}.
\bibitem{HMM24} I. Hu, D. van Melkebeek, and A. Morgan.
\emph{Polynomial Identity Testing via Evaluation of Rational Functions}.
Theory of Computing 20(1), 2024.
\url{https://doi.org/10.4086/toc.2024.v020a001}.
\bibitem{SS26} S. Singh and V. Singh.
\emph{An Algebraic Rigidity Framework for Order-Oblivious Deterministic
Black-Box PIT of ROABPs}. arXiv:2602.13449v1, 13 February 2026.
\url{https://arxiv.org/abs/2602.13449}.
\end{thebibliography}
\end{document}
